\scanpage{27}
\mspage{\handwritten{21}}
finie, \inserted{\handwritten{et à fibres affines,}} $Z$ un sous-préschéma en groupes. Pour que $Z$ soit un centre réductif de $G$, il faut et il suffit qu'il soit de type multiplicatif, et que pour tout $s\handwritten{\in}S$, \deleted{désignant par} $Z_s$ soit \deleted{\unsure{le}}un centre réductif de $G_s$.

En effet, il résulte d'abord de Exp IX 5.6. b\deleted{\unsure{1}}) qu'alors $Z$ est central. Comme les conditions envisagées sont stables par changement de base, il reste à prouver que tout homomorphisme central $u:H\handwritten{\to}G$, avec $H$ de type multiplicatif et de type fini, se factorise par $Z$. Or , $Z$ étant central, $u$ et l'immersion canonique $v:Z\handwritten{\to}G$ définissent un homomorphisme de groupes
\[
w:H\times_S Z\handwritten{\longrightarrow}G\handwritten{.}
\]
\handwritten{E}n vertu de Exp IX 6.8. ce dernier admet un groupe image $K$, qui est un sous-groupe de type multiplicatif de $G$, et tout revient à prouver que l'on a $K=Z$. Or cela est vrai fibre par fibre d'après l'hypothèse sur $Z$, et il suffit maintenant d'appliquer Exp IX 5.1. bi\handwritten{s,} \inserted{ce qui achève de prouver \handwritten{4}.3.}

On notera que dans le critère \handwritten{4}.3., l'hypothèse que pour tout $s\handwritten{\in}S$, $Z_s$ soit \deleted{\unsure{le}}un centre réductif de $G_s$ est en fait purement géométrique, i.e. il suffit de la vérifier sur la clôture algébrique de $k(s)$, comme il résulte de la deuxième assertion \handwritten{4}.2.

\margin{\handwritten{\unsure{$\nearrow$|1}}}
\uline{Théorème} \handwritten{4}.4. Soit $G$ un groupe algébrique \uline{affine} sur un corps $k$. Alors $G$ admet un centre réductif.

Comme le centre de $G$ est représentable par un sous-groupe fermé de $G$ (Exp VIII 6.7.) , on est ramené aussitôt au cas où $G$ est commuta\deleted{\unsure{ti|tif}}tif. De plus, en vertu de \handwritten{4}.2. on peut supposer $k$ algébriquement clos. Nous verrons \deleted{\unsure{plus loin|plus bas}}\inserted{\handwritten{dans Exp $\overline{\underline{\mathrm{XIII}}}$}} que dans ce cas $G$ s'écrit comme un produit \deleted{\unsure{et|de}} $Z\times U$, où $Z$ est de type multiplicatif, et $U$ \textquotedblleft unipotent'', d'où il résultera aussitôt que $Z$ est bien un centre réductif de $G$. Nous allons
% archive sheet number: 26
